\section{Мотивация и идея}

Больше всего вопросов при изучении современных средств
интерактивных доказательств (Lean, Rocq, Agda...)
вызывает устройство \textit{иерархий вселенных}:
$\mathtt{Type}_1 : \mathtt{Type}_2 : \ldots$.
Вызывают они вопросы и с метатеоретической точки зрения
--- в доказательствах здравости $\mathtt{Type}_n$
интерпретируется как $n$-й недостижимый кардинал.

Поэтому интересно искать альтернативные варианты
устройства иерархии. Далее описан забавный вариант,
располагающий типы по ``вселенным'' соответственно
\textit{рангу функций}. Термы в этой системе сильно
нормализуются, так что она может быть использована в
качестве основы для средства автоматической проверки
доказательств. Также нами была построена
теоретико-множественная модель этой системы в ZFC,
открывающая возможные применения нашей системы в высшей
теории вычислимости.

Б\'{о}льшая часть описанных результатов формально
верифицирована на Rocq.

\section{Синтаксис чистых систем типов}

Наша система является конкретным экземпляром
\textit{Чистой системы типов}. ЧСТ определяется поверх
множества \textit{термов}, параметризованного
множеством переменных $X$ и множеством сортов $L$:

\begin{figure}[h]
\begin{bnf}
$x$ : \textsf{Переменная} :in: $X$;;
$\ell$ : \textsf{Сорт} :in: $L$;;
$T$, $U$, $t$, $u$ : \textsf{Терм} ::=
| $x$ : переменная
| $\ell$ : сорт как тип
| $\Pi(x\,\colon T).\,U$ : тип зависимых функций
| $\lambda(x\,\colon T).\,u$ : анонимная функция
| $u\,\app\,t$ : применение функции
;;
\end{bnf}
\end{figure}

В формализации синтаксиса с операторами, связывающими
переменные, такими как $\Pi$ и $\lambda$, есть много
тонкостей, зависящих от выбора представления. В нашей
формализации мы выбрали представление синтаксиса как
\textit{относительной монады} --- тип термов
параметризован количеством свободных переменных;
применение операторов связывания увеличивает количество
свободных переменных в подтерме на $1$; переменные
представляют собой значения типа \textit{ограниченных
натуральных чисел}.

\begin{definition}[Ограниченные натуральные числа]
Тип $\Fin(n)$ состоит из натуральных чисел, строго
меньших $n$, с конструкторами
$\mathsf{fzero}:\Fin(n+1)$ (первый элемент) и
$\mathsf{fsucc}:\Fin(n)\to\Fin(n+1)$
(преемник элемента).
\end{definition}

Наш выбор представления мотивирован тем, что,
во-первых, в реализации избегающей связывания
подстановки для такой системы ошибиться гораздо
сложнее, чем для системы с наивным строковым
представлением переменных; во-вторых, наличие такого
параметра увеличивает количество информации, доступной
для элаборации движком Rocq, что уменьшает количество
аннотаций в формулировках теорем и в доказательствах
по сравнению с непараметризованными представлениями;
в-третьих, параметризация натуральным числом вместо
параметризации типом переменных позволяет естественным
образом связать термы и контекст в правиле типизации.

Далее описаны определения синтаксических операций над
термами, вспомогательных функций и связанных с ними
лемм, иллюстрирующих корректность определений.

\begin{definition}[Ослабление функции переименования]
Для функции переименования $f : \Fin(m) \to \Fin(n)$
\emph{ослабление} $\weak(f) : \Fin(m+1) \to \Fin(n+1)$
отображает новый индекс $\mathsf{fzero}$ в
$\mathsf{fzero}$ и сдвигает остальные согласно $f$.
\end{definition}

\begin{definition}[Переименование]
Для функции переименования $f : \Fin(m) \to \Fin(n)$
операция \emph{переименования} $\rename$ преобразует
терм в контексте длины $m$ в терм в контексте длины
$n$: вселенные остаются без изменений; в типе зависимых
функций область определения переименовывается по $f$,
а область значений --- по $\weak(f)$; индексы
переменных отображаются по $f$; в лямбде область
определения переименовывается по $f$, а тело --- по
$\weak(f)$; в применении обе части переименовываются
по $f$.
\end{definition}

\begin{definition}[Сдвиг]
\emph{Сдвиг} $\shift$ терма $t$ в контексте длины $n$
получает терм в контексте длины $n+1$ путём увеличения
каждого индекса переменной (т.\,е.~переименование по
функции, отображающей каждый индекс в его преемник).
\end{definition}

\begin{definition}[Транспонирование]
Для функции подстановки $f : \Fin(m) \to \Term(L, n)$
\emph{транспонирование} $\transpose(f) : \Fin(m+1) \to
\Term(L, n+1)$ по индексу $i$ возвращает терм в
контексте длины $n+1$: при $i = \mathsf{fzero}$ это
переменная $\var(\mathsf{fzero})$; при
$i = \mathsf{fsucc}(j)$ это $\shift(f(j))$.
\end{definition}

\begin{definition}[Подстановка]
Для функции подстановки $f : \Fin(m) \to \Term(L, n)$ и
терма $t$ в контексте длины $m$ операция
\emph{подстановки} $\replace(f, t)$ заменяет переменные
в $t$ соответствующими термами из $f$: вселенные
остаются без изменений; в типе зависимых функций
подстановка в область определения происходит по $f$, а
в область значений --- по транспонированию $f$;
переменная $i$ заменяется на $f(i)$; в лямбде
подстановка в область определения происходит по $f$, а
в тело --- по транспонированию $f$; в применении в обе 
части подставляется $f$.
\end{definition}

\begin{definition}[Подстановка в свежую переменную]
Подстановка терма $v:\Term(L,n)$ в свежую переменную
терма $t:\Term(L,n+1)$ (обозначение в Rocq:
$t\;\sub\;v$) есть операция замены переменной с
индексом $\mathsf{fzero}$ на $v$ и одновременного
сдвига остальных переменных на $-1$.
\end{definition}
В неформальном описании в тексте подстановка в свежую
переменную $x$ обозначается как $t[x \coloneqq v]$ для
читаемости.

\begin{definition}[Толкание]
Для терма $t:\Term(L,n)$ операция \emph{толкания}
строит функцию подстановки
$\mathsf{push}(t):\Fin(m+n+1)\to\Term(L,m+n)$, где
первые $m$ переменных остаются без изменений, $(m+1)$-я
отображается в $t$, а оставшиеся $n$ переменных
сдвигаются на $-1$.
\end{definition}

\begin{definition}[Индексная подстановка]
Подстановка терма $v$ в переменную с заданным индексом
в терме $t$ (обозначение $t\;\sub^i\;v$) выполняется
путём толкания $v$ к нужному индексу и последующего
применения полученной подстановки.
\end{definition}

\begin{lemma}[Экстенсиональность переименования]
Если две функции переименования $f$ и $g$ совпадают на
всех индексах, то они дают одинаковый результат при
применении к любому терму.
\end{lemma}

\begin{lemma}[Экстенсиональность транспонирования]
Если две функции $f$ и $g$ совпадают на всех индексах,
то их транспонирования совпадают на всех индексах.
\end{lemma}

\begin{lemma}[Экстенсиональность подстановки]
Если две функции подстановки $f$ и $g$ совпадают на
всех индексах, то они дают одинаковый результат при
применении к любому терму.
\end{lemma}

\begin{lemma}[Композиция переименований]
Переименование по $f$ после переименования по $g$
равносильно переименованию по композиции $f \circ g$.
\end{lemma}

\begin{lemma}
    [Переименование коммутирует с подстановкой]
Переименование по $f$ терма, к которому применялась
подстановка $g$, эквивалентно подстановке $g$,
переименованной по $f$.
\end{lemma}

\begin{lemma}
    [Подстановка коммутирует с переименованием]
Подстановка в переименованный терм равносильна
подстановке по композиции соответствующих функций.
\end{lemma}

\begin{lemma}[Подстановка с тождественным отображением]
Замена переменных на самих себя есть тождественная
операция на термах.
\end{lemma}

\begin{lemma}[Композиция подстановок]
Подстановка одной подстановки в другую равносильна
подстановке по композиции функций.
\end{lemma}

\begin{lemma}
    [Переименование распределяется над подстановкой]
Переименование подстановки $t\;\sub\;u$ по $f$
совпадает с подстановкой $\weak(f)(t)\;\sub\;f(u)$.
\end{lemma}

\begin{lemma}
    [Подстановка распределяется над подстановкой]
Подстановка в подстановку $t\;\sub\;u$ по $f$
совпадает с подстановкой
$\transpose(f)(t)\;\sub\;f(u)$.
\end{lemma}

\begin{lemma}[Подстановка в сдвинутый терм]
Подстановка в свежую переменную сдвинутого терма
восстанавливает исходный терм:
$\shift(t)\;\sub\;u = t$.
\end{lemma}

\begin{lemma}[Сдвиг коммутирует с индексной подстановкой]
Сдвиг, за которым следует индексная подстановка на
$m+1$, равносилен индексной подстановке на $m$, за
которой следует сдвиг.
\end{lemma}

\section{ЧСТ как система подстановок}

\subsection{Общие факты об отношениях}

Вспомним некоторые базовые определения из теории
подстановочных систем.

\begin{definition}[Рефлексивно-транзитивное замыкание]
Для множества $A$ и отношения $R$ на $A$
\emph{рефлексивно-транзитивное замыкание} $\RTC(R)$
есть наименьшее отношение на $A$, генерируемое двумя
конструкторами: конструктор рефлексивности, дающий
$x\;\RTC\;x$ для каждого $x \in A$; и конструктор шага,
который, по заданным $x\;R\;y$ и $y\;\RTC(R)\;z$, даёт
$x\;\RTC(R)\;z$.
\end{definition}

\begin{definition}[Симметрическое замыкание]
\emph{Симметрическое замыкание} $\Sym(R)$ есть
наименьшее отношение на $A$, генерируемое:
конструктором $\mathsf{forward}$, который по заданному
$x\;R\;y$ даёт $x\;\Sym(R)\;y$; и конструктором
$\mathsf{backward}$, который по заданному $y\;R\;x$
даёт $x\;\Sym(R)\;y$.
\end{definition}

\begin{definition}[Обратное отношение]
\emph{Обратное} отношение $R^{-1}$ определяется так: \\
$x\;R^{-1}\;y$ тогда и только тогда, когда $y\;R\;x$.
\end{definition}

\begin{definition}[Замыкание до эквивалентности]
\emph{Замыкание отношения $R$ до эквивалентности} есть 
рефлексивно-транзитивное замыкание симметрического
замыкания:\\
$\EquivClosure R \;\triangleq\; \RTC(\Sym R)$.
\end{definition}

\begin{definition}[Прообраз отношения]
Для функции $f : A \to B$ и отношения $R$ на $B$ \emph{прообраз}
отношения $R$ по $f$ (обозначение $R \mathbin{@} f$) есть отношение на $A$,
определённое тем, что
$x\;(R \mathbin{@} f)\;y$ тогда и только тогда, когда $f(x)\;R\;f(y)$.
\end{definition}

\begin{definition}[Свойство Чёрча--Россера]
Отношение $R$ на $A$ обладает
\emph{свойством Чёрча--Россера}, если для всех
$t, t_1, t_2 \in A$ таких, что $t\;R\;t_1$ и
$t\;R\;t_2$, существует $t^*$ такой, что $t_1\;R\;t^*$
и $t_2\;R\;t^*$:
\[\begin{tikzcd}
    t \arrow[r] \arrow[d]
    & t_1 \arrow[d, dashed] \\
    t_2 \arrow[r, dashed] & t^*
\end{tikzcd}\]
\end{definition}

Далее, словом \textbf{Экземпляр} обозначены факты,
которые были сохранены нами в базе экземпляров Rocq для
автоматического вывода свойств конкретных отношений,
введённых нами в ходе формализации.

\begin{instance}
Рефлексивно-транзитивное замыкание любого отношения
рефлексивно.
\end{instance}

\begin{instance}
Рефлексивно-транзитивное замыкание любого отношения
транзитивно.
\end{instance}

\begin{lemma}
    Если $x\;R\;y$, то $x\;\RTC(R)\;y$.
\end{lemma}

\begin{instance}
Если $R$ симметрично, то его рефлексивно-транзитивное
замыкание также симметрично.
\end{instance}

\begin{instance}
Симметрическое замыкание любого отношения симметрично.
\end{instance}

\begin{instance}
Замыкание любого отношения до эквивалентности есть
отношение эквивалентности
(т.\,е.~является рефлексивным, симметричным и
транзитивным).
\end{instance}

\begin{lemma}
Если $x\;\RTC(R^{-1})\;y$, то $y\;\RTC(R)\;x$.
\end{lemma}

\begin{lemma}[
    Отображение через рефлексивно-транзитивное
    замыкание
]
Для функции $f : A \to B$ и отношения $R$ на $B$: если
$x\;\RTC(R \mathbin{@} f)\;y$, то $f(x)\;\RTC(R)\;f(y)$.
\end{lemma}

\begin{lemma}
    [Монотонность рефлексивно-транзитивного замыкания]
Если $Q \subseteq R$ (т.\,е.~$Q\,x\,y \Rightarrow R\,x\,y$ для всех
$x, y$), то $\RTC Q \subseteq \RTC R$.
\end{lemma}

\begin{lemma}
    [Отображение через симметрическое замыкание]
Для функции $f : A \to B$ и отношения $R$ на $B$: если
$x\;\Sym(R \mathbin{@} f)\;y$, то $f(x)\;\Sym(R)\;f(y)$.
\end{lemma}

\begin{lemma}[Монотонность симметрического замыкания]
Если $Q \subseteq R$, то $\Sym Q \subseteq \Sym R$.
\end{lemma}

\begin{lemma}
Если $x\;R\;y$, то $x\;\EquivClosure(R)\;y$.
\end{lemma}

\begin{lemma}
    [Отображение через замыкание до эквивалентности]
Для функции $f : A \to B$ и отношения $R$ на $B$: если
$x\;\EquivClosure(R \mathbin{@} f)\;y$, то
$f(x)\;\EquivClosure(R)\;f(y)$.
\end{lemma}

\begin{lemma}
    [Монотонность замыкания до эквивалентности]
Если $Q \subseteq R$, то \\
$\EquivClosure(Q) \subseteq \EquivClosure(R)$.
\end{lemma}

\begin{lemma}
    [Экстенсиональность свойства Чёрча--Россера]
Если два отношения $Q$ и $R$ на $A$ равны
(т.\,е.~$Q(a,b) \Leftrightarrow R(a,b)$ для всех $a, b$), то $Q$
обладает свойством Чёрча--Россера тогда и только тогда,
когда $R$ обладает свойством Чёрча--Россера.
\end{lemma}

\begin{theorem}[Свойство Чёрча--Россера для
    рефлексивно-транзитивного замыкания]
Если отношение $R$ обладает свойством Чёрча--Россера,
то его рефлексивно-транзитивное замыкание $\RTC R$
также обладает свойством Чёрча--Россера.
\end{theorem}

\subsection{Конгруэнции}

В этом разделе мы определим, что такое конгруэнция на
термах. Здесь и далее, словом \textbf{Класс} мы будем
обозначать определения, соответствие которым удобно
считать \textbf{Экземплярами} для автоматического
вывода свойств в Rocq.

\begin{class}[Конгруэнция]
Семейство отношений $R$ на термах (индексируемое длиной
контекста) является \emph{конгруэнцией}, если для него
выполнен набор из шести условий совместимости:
\begin{itemize}
\item отношение совместимо с типом функций по области
    определения (если $T \mathrel{R} T'$, то
    $(\PiT(x:T).U) \mathrel{R} (\PiT(x:T').U)$) и по
    области значений (если $U \mathrel{R} U'$, то
    $(\PiT(x:T).U) \mathrel{R} (\PiT(x:T).U')$);
\item отношение совместимо с лямбдами по области
    определения (если $T \mathrel{R} T'$, то \\
    $(\lam(x:T).t) \mathrel{R} (\lam(x:T').t)$) и по
    телу (если $t \mathrel{R} t'$, то
    $(\lam(x:T).t) \mathrel{R} (\lam(x:T).t')$);
\item отношение совместимо с аппликациями по функции
    (если $t \mathrel{R} t'$, то
    $(t\;\app\;u) \mathrel{R} (t'\;\app\;u)$) и по
    аргументу (если $u \mathrel{R} u'$, то
    $(t\;\app\;u) \mathrel{R} (t\;\app\;u')$).
\end{itemize}
\end{class}

\begin{instance}
    [Симметрическое замыкание сохраняет конгруэнтность]
Если $R$ --- конгруэнция, то семейство симметрических
замыканий $\Sym(R)$ также является конгруэнцией.
\end{instance}

\begin{instance}
    [Рефлексивно-транзитивное замыкание сохраняет
    конгруэнтность]
Если $R$ --- конгруэнция, то семейство
рефлексивно-транзитивных замыканий $\RTC(R)$ также
является конгруэнцией.
\end{instance}

\begin{lemma}
    [Параллельная конгруэнтность для типов функций]
Если $R$ --- транзитивная конгруэнция, а
$T \mathrel{R} T'$ и $U \mathrel{R} U'$, то
$(\PiT(x:T).U) \mathrel{R} (\PiT(x:T').U')$.
\end{lemma}

\begin{lemma}[Параллельная конгруэнтность для лямбд]
Если $R$ --- транзитивная конгруэнция, а
$T \mathrel{R} T'$ и $t \mathrel{R} t'$, то
$(\lam(x:T).t) \mathrel{R} (\lam(x:T').t')$.
\end{lemma}

\begin{lemma}
    [Параллельная конгруэнтность для аппликаций]
Если $R$ --- транзитивная конгруэнция, а
$t \mathrel{R} t'$ и $v \mathrel{R} v'$, то
$(t\;\app\;v) \mathrel{R} (t'\;\app\;v')$.
\end{lemma}

\subsection{Отношение одношаговой редукции}

Наконец, мы готовы вернуться к описанию чистой системы
типов. Следующий шаг --- определение редукций, т.е.
вычислений в системе. Одношаговая редукция термов
определяется как наименьшее конгруэнтное отношение,
выполняющее правило бета-редукции о применении
анонимной функции к аргументу:

\begin{mathpar}
\inferrule{T \to_\beta T'}
    {\Pi(x:T).U \to_\beta \Pi(x:T').U}
\and
\inferrule{U \to_\beta U'}
    {\Pi(x:T).U \to_\beta \Pi(x:T).U'}
\and
\inferrule{T \to_\beta T'}
    {\lambda(x:T).t \to_\beta \lambda(x:T').t}
\and
\inferrule{t \to_\beta t'}
    {\lambda(x:T).t \to_\beta \lambda(x:T).t'}
\and
\inferrule{t \to_\beta t'}
    {t\,\app\,u \to_\beta t'\,\app\,u}
\and
\inferrule{u \to_\beta u'}
    {t\,\app\,u \to_\beta t\,\app\,u'}
\and
\inferrule{}
    {(\lambda(x:T).t)\,\app\,u \to_\beta
    t[x\coloneqq u]}
\end{mathpar}

Ниже --- несколько лемм, показывающих взаимоотношения
между редукцией и уже определёнными операциями на
термах.

\begin{instance}
Одношаговая $\beta$-редукция является конгруэнцией.
\end{instance}

\begin{definition}[Многошаговая $\beta$-редукция]
Отношение $\mstep$ есть \\ рефлексивно-транзитивное
замыкание отношения $\bstep$.
\end{definition}

\begin{lemma}[Переименование коммутирует с шагом]
Если $t \bstep t'$, то
$\rename(f)(t) \bstep \rename(f)(t')$.
\end{lemma}

\begin{lemma}[Подстановка коммутирует с шагом]
Если $t \bstep t'$, то
$\replace(f)(t) \bstep \replace(f)(t')$.
\end{lemma}

\begin{lemma}[Лямбда-инверсия для переименования]
Если переименование терма совпадает с некоторой
анонимной функцией, то исходный терм сам
является анонимной функцией, и переименование
раскладывается соответственно.
\end{lemma}

\begin{lemma}[Инверсия шага для переименования]
Если $\rename(f)(t) \bstep u$, то существует $v$ такой,
что $u = \rename(f)(v)$ и $t \bstep v$.
\end{lemma}

\begin{lemma}
    [Переименование коммутирует с многошаговой
    редукцией]
Если $t \mstep t'$, то
$\rename(f)(t) \mstep \rename(f)(t')$.
\end{lemma}

\begin{lemma}
    [Транспонирование коммутирует с многошаговой
    редукцией]
Если каждый компонент функции подстановки преобразуется
в соответствующий компонент другой, то их
транспонирования также преобразуются друг в друга.
\end{lemma}

\begin{lemma}
    [Подстановка коммутирует с многошаговой редукцией]
Если каждый компонент функции подстановки преобразуется
в соответствующий компонент другой, то подставленные
термы также преобразуются друг в друга.
\end{lemma}

\begin{lemma}[Представление произведения]
Если $\PiT(x:T).U \mstep t$, то $t$ является типом
зависимых функций.
\end{lemma}

\begin{lemma}
    [Инверсия произведения для многошаговой редукции]
\.\\ Если $\PiT(x:T).U \mstep \PiT(x:T').U'$, то
$T \mstep T'$ и $U \mstep U'$.
\end{lemma}

\subsection{Конфлюэнтность редукций}

Классическим результатом в теории ЧСТ является
конфлюэнтность редукций, т.е. выполнение свойства
Чёрча-Россера для многошаговой редукции. В этой секции
мы воспроизводим это доказательство на основе
доказательства для нетипизированного лямбда-исчисления
\cite{CRproof}. Ключевой элемент доказательства ---
дополнительное построение \emph{параллельной
$\beta$-редукции}, производящей все доступные в терме
подстановки ``одновременно'' и обладающей свойством
Чёрча-Россера.

\begin{definition}[Параллельная $\beta$-редукция]
\begin{mathpar}
\inferrule{ }{t \pstep t}
\and
\inferrule{T \pstep T' \\ U \pstep U'}
    {\Pi(x:T).U \pstep \Pi(x:T').U'}
\and
\inferrule{T \pstep T' \\ t \pstep t'}
    {\lambda(x:T).t \pstep \lambda(x:T').t'}
\\
\inferrule{t \pstep t' \\ u \pstep u'}
    {t\,\app\,u \pstep t'\,\app\,u'}
\and
\inferrule{t \pstep t' \\ u \pstep u'}
    {(\lambda(x:T).t)\,\app\,u \pstep
    t'[x\coloneqq u']}
\end{mathpar}
\end{definition}

\begin{lemma}
    Если $t \bstep t'$, то $t \pstep t'$.
\end{lemma}

\begin{lemma}
    Если $t \pstep t'$, то $t \mstep t'$.
\end{lemma}

\begin{lemma}
Если $t \pstep t'$, то
$\rename(f)(t) \pstep \rename(f)(t')$.
\end{lemma}

\begin{lemma}
    [Транспонирование коммутирует с параллельным шагом]
Если каждый компонент функции подстановки делает
параллельный шаг в соответствующий компонент другой, то
их транспонирования также делают параллельный шаг.
\end{lemma}

\begin{lemma}
    [Подстановка коммутирует с параллельным шагом]
Если каждый компонент функции подстановки делает
параллельный шаг в соответствующий компонент другой, и
$t \pstep t'$, то подставленные термы также делают
параллельный шаг друг в друга.
\end{lemma}

\begin{lemma}
Если $t \pstep t'$, то $f\;\sub\;t \pstep f\;\sub\;t'$.
\end{lemma}

\begin{lemma}
Если $f \pstep f'$ и $t \pstep t'$, то
$f\;\sub\;t \pstep f'\;\sub\;t'$.
\end{lemma}

\begin{lemma}[Лямбда-инверсия для параллельного шага]
Если анонимная функция делает параллельный шаг в
некоторый терм $t'$, то $t'$ является лямбдой, и
параллельный шаг раскладывается на параллельные шаги в
области определения и теле лямбды.
\end{lemma}

\begin{lemma}
    [Конфлюэнтность параллельной редукции]
Отношение параллельной $\beta$-редукции обладает
свойством Чёрча--Россера.
\end{lemma}

\begin{theorem}
    [Конфлюэнтность многошаговой редукции]
Отношение многошаговой $\beta$-редукции обладает
свойством Чёрча--Россера.
\end{theorem}

\subsection{Вычислительное равенство термов}

Ключевым понятием в определении правил типизации ЧСТ
является \textit{вычислительное равенство} термов.

\begin{definition}[Вычислительное равенство]
Два терма $t$ и $t'$ \emph{равны вычислительно}
(обозначение $t \Deq t'$), если они связаны замыканием
одношагового $\beta$-преобразования до эквивалентности.
\end{definition}

Естественным образом, оно уважается базовыми
синтаксическими операциями.

\begin{lemma}
Если $t \Deq t'$, то
$\rename(f)(t) \Deq \rename(f)(t')$.
\end{lemma}

\begin{lemma}
Если $t \Deq t'$, то
$\replace(f)(t) \Deq \replace(f)(t')$.
\end{lemma}

Кроме этого, благодаря конфлюэнтности системы,
вычислительное равенство обладает рядом важных свойств.

\begin{theorem}
    [Характеризация вычислительного равенства]
Два терма равны вычислительно тогда и только тогда,
когда существует их общий результат редукции:
т.\,е.~$t \Deq t'$ тогда и только тогда, когда
существует терм $t^*$ такой, что $t \mstep t^*$ и
$t' \mstep t^*$.
\end{theorem}

\begin{lemma}
    [Инверсия типа функций для вычислительного
    равенства]
Если $\PiT(x:T).U \Deq \PiT(x:T').U'$, то
$T \Deq T'$ и $U \Deq U'$.
\end{lemma}

\begin{lemma}
    [Подстановка сохраняет вычислительное равенство]
Если $t \Deq t'$, то
$f\;\sub\;t \Deq f\;\sub\;t'$.
\end{lemma}

\section{Типизация ЧСТ}

\subsection{Контекст}

Типизация языков с операторами связывания переменных не
обходится без \textit{контекста} --- структуры данных,
хранящей тип каждой доступной переменной. В случае
чистых систем типов, множество типов и термов
совпадает, так что это просто список термов,
выступающих в роли типа; тонкость в том, что каждый
следующий тип в контексте может зависеть от всех ранее
введённых переменных.

\begin{definition}[Контексты]
\emph{Контекст} длины $n$ --- это либо пустой контекст
$\cdot$ (длины $0$), либо расширение $\Gamma, x:T$
контекста $\Gamma$ длины $n$ типом (термом) $T$ в
контексте длины $n$, дающее контекст длины $n+1$.
\end{definition}

\begin{definition}[Предшественник контекста]
Для непустого контекста $\Gamma = \Gamma_0,x:T$
\emph{предшественник} есть $\Gamma_0$.
\end{definition}

\begin{definition}[Вершина контекста]
Для непустого контекста $\Gamma = \Gamma_0,x:T$
\emph{вершина} (или головной элемент) есть $T$.
\end{definition}

\begin{definition}[Индексирование контекста]
Для контекста $\Gamma$ длины $n$ и индекса
$i \in \Fin(n)$ (обозначение: $\Gamma \mathbin{!!} i$)
извлекает $i$-й тип из контекста (считая с вершины),
с надлежащим сдвигом.
\end{definition}

В неформальном изложении далее индексы отождествляются
с переменными, так что допустимо обозначение
$\Gamma\mathbin{!!}x$.

\begin{definition}[Переименователь ослабления]
Функция $\wrr : \Fin(m+n) \to \Fin(m+n+1)$,
сдвигающая индексы с учётом нового типа, вставленного
на позицию $m$, считая с вершины контекста.
\end{definition}

\begin{definition}[Вставка в контекст]
Вставка типа $T:\Term(L,n)$ в контекст $\Gamma$ длины
$m+n$ (на позицию $m$) даёт контекст длины $m+n+1$
с надлежащим переименованием существующих типов.
\end{definition}

\begin{definition}[Сжатие контекста]
Удаление первых $m$ типов из контекста длины $m+n$,
оставляя контекст длины $n$ ($\shrink$).
\end{definition}

\begin{definition}[Индексная подстановка в контекст]
Подстановка терма $t$ в переменную с заданной позицией
в контексте, дающая новый контекст, в котором эта
переменная заменена на $t$ с надлежащим индексом.
\end{definition}

\begin{definition}[Сжатие до одного элемента]
Извлечение ($\squeeze$) типа $(m+1)$-й переменной из
контекста длины $m+n+1$ путём снятия первых $m$ типов.
\end{definition}

\begin{lemma}
Применение переименователя ослабления к преемнику
индекса равно преемнику применения переименователя
ослабления к индексу.
\end{lemma}

\begin{lemma}
Ослабление переименователя ослабления равно
самому переименователю ослабления на преемниках
индексов.
\end{lemma}

\begin{lemma}
Переименование по ослаблению переименователя ослабления
равносильно переименованию по переименователю
ослабления.
\end{lemma}

\begin{lemma}[Преемник индекса]
Индексирование контекста, расширенного новым типом,
по преемнику индекса возвращает сдвинутый результат
индексирования индексом в исходном контексте.
\end{lemma}

\begin{lemma}[Вставка и индексирование]
Индексирование вставленного контекста по ослабленному
индексу даёт переименование исходного индекса.
\end{lemma}

\subsection{Сорт}

Последний и самый важный компонент, необходимый для
введения правил типизации --- сигнатура сортов,
задающая взаимное отношение сортов друг между другом и
отнесение типов зависимых функций к тем или иным
сортам.

\begin{class}[Сигнатура сортов]
\emph{Сигнатура сортов} на типе сортов $L$ состоит из
бинарного отношения $\ax$ ($\ax(\ell,\ell')$
постулирует, что терм $\ell$ имеет тип $\ell'$) и
трёхместного отношения $\ruleL$
($\ruleL(\ell_S,\ell_T,\ell)$ постулирует, что если $S$
--- тип сорта $\ell_S$, а $T$ --- тип сорта $\ell_T$,
то $\Pi(x:S).T$ --- тип сорта $\ell$).
\end{class}

\begin{class}[Функциональная сигнатура сортов]
Сигнатура сортов \emph{функциональна}, если:
отношение $\ax$ функционально (если
$\ax(\ell, \ell_1)$ и $\ax(\ell, \ell_2)$, то
$\ell_1 = \ell_2$); и отношение $\ruleL$
функционально по первой паре аргументов (если
$\ruleL(\ell_t, \ell_u, \ell_1)$ и
$\ruleL(\ell_t, \ell_u, \ell_2)$, то
$\ell_1 = \ell_2$).
\end{class}

\begin{class}[Полная сигнатура сортов]
Сигнатура сортов \emph{полна}, если:
каждый сорт находится в каком-то другом (для каждого
$\ell$ существует $\ell'$ такой, что
$\ax(\ell, \ell')$); и каждая пара сортов определяет
сорт соответствующих типов функций (для каждого
$\ell_t, \ell_u$ существует $\ell$ такой,
что $\ruleL(\ell_t, \ell_u, \ell)$).
\end{class}

\subsection{Правила типизации}

Здесь и далее фиксируется некоторая произвольная
сигнатура сортов.

\begin{definition}[Суждение типизации]
\begin{mathpar}
\inferrule{ }{\Gamma\vdash x:\Gamma\mathbin{!!}x}
\and
\inferrule{ }{\Gamma\vdash s:s'} (s, s')\in\mathrm{Ax}
\and
\inferrule{
    \Gamma\vdash T:s_T \\
    \Gamma,x:T\vdash U:s_U
}{\Gamma\vdash\Pi(x:T).\,U:s}
(s_T, s_U, s)\in\mathrm{Rule}
\and
\inferrule{
    \Gamma\vdash T:s_T \\
    \Gamma,x:T\vdash u:U
}{\Gamma\vdash\lambda(x:T).u:\Pi(x:T).U}
\and
\inferrule{
    \Gamma\vdash u:\Pi(x:T).U \\
    \Gamma\vdash t:T
}{\Gamma\vdash u\,\app\,t:U[x\coloneqq t]}
\and
\inferrule{
    \Gamma\vdash t:T \\
    \Gamma\vdash U:s \\
    T =_\beta U
}{\Gamma\vdash t:U}
\end{mathpar}
\end{definition}

Благодаря введённым ранее определениям для суждения
типизации можно сформулировать и доказать усиленные
версии следующих структурных свойств, из которых
следуют их обычные версии, необходимые для
доказательства сохранения типизации.

\begin{lemma}[Сильное ослабление]
Если $\Gamma \vdash t : T$ и мы вставляем новый
тип $U$ в контекст, то переименованный терм
и тип остаются корректно типированными: \\
$\insertt(\Gamma, U) \vdash \rename(\wrr)(t) : \rename(\wrr)(T)$.
\end{lemma}

\begin{lemma}[Ослабление]
Если $\Gamma \vdash t : T$, то для любого типа $U$
верно \\ $\Gamma,x:U \vdash \shift(t) : \shift(T)$.
\end{lemma}

\begin{lemma}[Сильная подстановка]
Если $\Gamma \vdash t : T$
и $\shrink(\Gamma) \vdash u : \squeeze(\Gamma)$,
то индексированная подстановка сохраняет суждение
типирования, т.е.
$\Gamma\;\sub^i\;u\vdash t\;\sub^i\;u:T\;\sub^i\;u$.
\end{lemma}

\begin{lemma}[Подстановка]
Если $\Gamma,x:U\vdash t:T$ и $\Gamma\vdash u:U$, то
$\Gamma\vdash t[x\coloneqq u]:T[x\coloneqq u]$.
\end{lemma}

\subsection{Леммы об инверсии}

Доказательство сохранения типизации существенно
опирается на существование \emph{лемм об инверсии},
выводящих типизацию подтермов из типизации терма.

\begin{lemma}[Инверсия сорта]
Если $\Gamma \vdash \ell : T$, то существует
сорт $\ell'$ такой, что $\ax(\ell, \ell')$ и
$T \Deq \ell'$.
\end{lemma}

\begin{lemma}[Инверсия типа функций]
Если $\Gamma \vdash \PiT(x:T).U : V$, то
существуют сорты $\ell_t, \ell_u, \ell$ такие, что
\[\left\{\begin{array}{l}
    \ruleL(\ell_t, \ell_u, \ell) \\
    \Gamma \vdash T : \ell_t \\
    \Gamma,x:T \vdash U : \ell_u \\
    V \Deq \ell
\end{array}\right.\]
\end{lemma}

\begin{lemma}[Инверсия переменной]
Если $\Gamma \vdash x : T$, то
$T \Deq \Gamma \mathbin{!!} x$.
\end{lemma}

\begin{lemma}[Инверсия лямбды]
Если $\Gamma \vdash \lam(x:T).t : V$, то
существуют сорт $\ell$ и терм $U$ такие, что
\[\left\{\begin{array}{l}
    \Gamma \vdash T : \ell \\
    \Gamma,x:T \vdash t : U \\
    V \Deq \PiT(x:T).U
\end{array}\right.\]
\end{lemma}

\begin{lemma}[Инверсия применения]
Если $\Gamma \vdash t\;\app\;u : V$, то
существуют термы $U$ и $T$ такие, что
\[\left\{\begin{array}{l}
    \Gamma \vdash t : \PiT(x:U).T \\
    \Gamma \vdash u : U \\
    V \Deq T[x\coloneqq u]
\end{array}\right.\]
\end{lemma}

\begin{theorem}
    [Однозначность типизации / Единственность типа]
Если сигнатура сортов функциональна,
$\Gamma \vdash t : T$ и
$\Gamma \vdash t : T'$, то $T \Deq T'$.
Иными словами, терм имеет не более одного типа
с точностью до вычислительного равенства.
\end{theorem}

\subsection{Сохранение типизации}

Мы подобрались к одному из самых важных свойств системы
--- редукция типизированного терма типизируется так же.
Как и с конфлюэнтностью, доказательство этого факта
требует некоторых дополнительных построений;
доказательство адаптировано из описания доказательства
непротиворечивости Rocq на Rocq (TODO: ссылка).

Важным элементом доказательства сохранения типизации
системы является определение одношаговой редукции на
контекстах.

\begin{definition}[$\beta$-редукция на контекстах]
\begin{mathpar}
\inferrule{T\bstep U}{\Gamma,x:T\dstep\Gamma,x:U}
\and\inferrule{\Gamma\dstep\Delta}
    {\Gamma,x:T\dstep\Delta,x:T}
\end{mathpar}
\end{definition}

Заметим, что в наших правилах типизации никак не
проверяется контекст; пришло время это сделать.

\begin{definition}[Корректные контексты]
Контекст $\Gamma$ \emph{корректен}
(обозначение $\wf(\Gamma)$) в одном из двух случаев:
\begin{enumerate}
    \item $\Gamma$ --- пустой;
    \item $\Gamma=\Gamma_0,x:T$; при этом $\Gamma$
        корректен и $\Gamma \vdash T : \ell$ для
        некоторого сорта $\ell$.
\end{enumerate}
\end{definition}

\begin{lemma}[Сильная редукция контекста]
Если $\Gamma \vdash t : T$ и $\Gamma$ корректен, а
также $\Gamma \dstep \Delta$, то
$\Delta \vdash t : T$.
Иными словами, редукция типа в контексте
сохраняет типизацию терма.
\end{lemma}

\begin{lemma}[Редукция контекста]
Если $\Gamma,x:U \vdash t : T$ и $\Gamma$ корректен,
а также $\Gamma \vdash U : \ell$ и $U \bstep U'$, то
$\Gamma,x:U' \vdash t : T$.
\end{lemma}

\begin{lemma}[Корректность типа]
Если сигнатура сортов полна, $\Gamma \vdash t : T$ и
$\Gamma$ корректен, то существует сорт $\ell$ такой,
что $\Gamma \vdash T : \ell$.
Иными словами, тип в суждении типизации тоже
типизируется некоторым сортом.
\end{lemma}

\begin{theorem}[Сохранение типизации]
Если сигнатура сортов полна, $\Gamma \vdash t : T$ и
$\Gamma$ корректен, а также $t \bstep t'$, то
$\Gamma \vdash t' : T$.
\end{theorem}

\section{О здравости чистых систем типов}

\subsection{Сильная нормализация}

Так, мы выяснили, что наша система не совсем
безнадёжна: вычисления не ``расходятся'', и типизация
сохраняется при вычислениях. Как видно из правил
типизации, для практической реализации этой системы в
виде средства интерактивных доказательств крайне
желательно уметь проверять вычислительное равенство
типов. Самый естественный способ --- по теореме о
характеризации: упрощаем термы по $\beta$-редукции
до тех пор, пока это возможно, и сравним результаты
синтаксически. Но вдруг найдётся такой терм, который
можно редуцировать бесконечно? Давайте исследуем
свойство \emph{нормализации}: терм нормализуется, если
любая цепочка его упрощений конечна. Формально это
определяется так:

\begin{definition}[Доступность]
Для отношения $R$ на множестве $A$ элемент $x$
\emph{доступен} (относительно $R$), если каждый
$y$ такой, что $x\;R\;y$, сам доступен.
\end{definition}

\begin{definition}[Сильная нормализуемость]
Терм $t$ \emph{сильно нормализуем}
(обозначение $\SN(t)$), если он доступен
относительно отношения одношаговой $\beta$-редукции.
\end{definition}

Какие свойства выполнены для сильно нормализующихся
термов?

\begin{lemma}[Сильная нормализуемость по голове]
Если $t\;\app\;u$ сильно нормализуем, то $t$
сильно нормализуем.
\end{lemma}

\begin{lemma}[Переименование сохраняет сильную нормализуемость]
Если терм сильно нормализуем, то любое его
переименование также сильно нормализуемо.
\end{lemma}

\begin{definition}[Аджаст]
Функция, которая для заданного $m$ отображает индекс
в $\Fin(n)$ в индекс в $\Fin(m+n)$ путём сдвига
всех индексов на $m$.
\end{definition}

\begin{lemma}[Аджаст и сильная нормализуемость]
Если терм, переименованный функцией $\adjust$,
сильно нормализуем, то исходный терм
сильно нормализуем.
\end{lemma}

\begin{definition}[Атомарные термы]
Терм \emph{атомарен}, если он есть: сорт; тип функций;
переменная; или применение атомарной функции к
атомарному аргументу.
\end{definition}

\begin{lemma}[Атомарность сохраняется при переименовании]
Если терм атомарен, то любое его переименование
также атомарно.
\end{lemma}

\begin{lemma}[Атомарность сохраняется при $\beta$-преобразовании]
Если $t \bstep t'$ и $t$ атомарен, то $t'$
также атомарен.
\end{lemma}

\begin{lemma}[Сильная нормализуемость атомарного применения]
Если $t$ атомарен и сильно нормализуем,
и $u$ сильно нормализуем, то $t\;\app\;u$
сильно нормализуем.
\end{lemma}

\begin{definition}[Слабое головное расширение]
Отношение $\WHE$ (слабое головное расширение)
на термах определяется так: если $t$ слабо головно
расширяется в $t'$, то $t\;\app\;u$ слабо головно
расширяется в $t'\;\app\;u$; и если $U$ и $u$
сильно нормализуемы, то
$\lam\,U\;t\;\app\;u$ слабо головно расширяется
в $t \sub u$.

Другими словами, $t$ слабо головно расширяется в $t'$,
если $t$ --- это цепочка применений аргументов к
сильно нормализуемому лямбда-выражению, где самый
первый первый аргумент также сильно нормализуем, а $t'$
--- терм, получающийся после бета-редукции в голове
этой цепочки.
\end{definition}

\begin{lemma}
    [Переименование коммутирует со слабым головным
    расширением]
Если $t$ слабо головно расширяется в $t'$,
то любое переименование $t$ слабо головно
расширяется в переименование $t'$.
\end{lemma}

\begin{lemma}
    [Слабое головное расширение
    и сильная нормализуемость]
Если $t$ слабо головно расширяется в $t'$
и $t'\;\app\;u$ сильно нормализуем, то
$t\;\app\;u$ сильно нормализуем.
\end{lemma}

В общем случае, проверка терма на нормализацию является
неразрешимой задачей, поскольку в текущем представлении
чистой системы типов очевидным образом выражается
Тьюринг-полное нетипизированное лямбда-исчисление.
Можно было бы надеяться, что типизация исправит
ситуацию, но она всё ещё автоматически не гарантирует
нормализацию, и какого-то классифицирующего критерия
для ЧСТ, позволяющего выяснить (не)завершимость термов
по сигнатуре сортов, нет. Тем не менее, для следующих
классических чистых систем типов была доказана сильная
нормализация:

\begin{itemize}
    \item Исчисление конструкций:
        \[\begin{array}{c}
            S=\{*,\square\}
            \quad \mathrm{Ax}=\{(*,\square)\}
            \\ \mathrm{Rule}=\{
                (*,*,*),(*,\square,\square),
                (\square,*,*),(\square,\square,\square)
            \}
        \end{array}\]
    \item Первая счётная иерархия \textit{вселенных}:
        \[\begin{array}{c}
            S=\{\mathcal{U}_\ell\;|\;\ell\in\mathbb{N}\}
            \quad \mathrm{Ax}=\{
                (\mathcal{U}_\ell,\mathcal{U}_{\ell+1})
                \;|\;\ell\in\mathbb{N}
            \}
            \\ \mathrm{Rule}=\{
                (\mathcal{U}_\ell, \mathcal{U}_{\ell'},
                \mathcal{U}_{\max(\ell,\ell')}) \;|\;
                \ell,\ell'\in\mathbb{N}
            \}
        \end{array}\]
    \item Счётная иерархия с импредикативным
        \texttt{Prop}:
        \[\begin{array}{rcl}
            S &=&\{\mathtt{Prop}\}\sqcup
            \{\mathcal{U}_\ell\;|\;\ell\in\mathbb{N}\}
            \\ \mathrm{Ax} &=&
            \{(\mathtt{Prop},\mathcal{U}_0)
            \}\sqcup\{
                (\mathcal{U}_\ell,\mathcal{U}_{\ell+1})
                \;|\;\ell\in\mathbb{N}
            \}
            \\ \mathrm{Rule} &=&\{
                (\mathcal{U}_\ell,\mathtt{Prop},
                \mathtt{Prop}) \;|\;
                \ell\in\mathbb{N}
            \} \;\sqcup
            \\ & \sqcup & \{
                (\mathtt{Prop}, s, s) \;|\; s\in S
            \} \;\sqcup
            \\ & \sqcup & \{
                (\mathcal{U}_\ell, \mathcal{U}_{\ell'},
                \mathcal{U}_{\max(\ell,\ell')}) \;|\;
                \ell,\ell'\in\mathbb{N}
            \}
        \end{array}\]
\end{itemize}

Кроме этого, о завершимости широкого класса чистых
систем типов есть следующий отрицательный результат.

\begin{theorem}[Girard'72]
Следующая ЧСТ (известная как System U) содержит
\textbf{не нормализующийся} терм:
\[\begin{array}{c}
    S=\{*,\square,\triangle\}
    \quad\mathrm{Ax}=\{
        (*,\square),(\square,\triangle)
    \}
    \\ \mathrm{Rule}=\{
        (*,*,*),(\square,*,*),(\triangle,*,*),
        (\square,\square,\square),
        (\triangle,\square,\square)
    \}
\end{array}\]
Более того, этот терм имеет тип $\Pi(T:*).T$.
\end{theorem}

Из не-завершимости термов в System U следует
не-завершимость любой другой системы, в которой
System U выражается. Другими словами, <<нормализующаяся
ЧСТ не может поддерживать одновременно \textbf{и}
импредикативный \texttt{Prop}, \textbf{и} полиморфные
функции>>. Например, из парадокса Жирара выводится
противоречивость ещё более простой ЧСТ, известной как
Type-in-Type:
\[S=\{*\}\quad\mathrm{Ax}=\{(*,*)\}\quad
\mathrm{Rule}=\{(*,*,*)\}\]

\subsection{Теоретико-множественная семантика}
Другое свойство, которое стоит
требовать от системы, претендующей на фундамент системы
доказательств --- каноничность: всякий терм вычисляется
до некоторого своего ``значения''. 
Кроме сильной нормализации, можно также исследовать и
другое свойство, наиболее важное для формализации
математики --- наличие теоретико-множественной модели.
У счётных иерархий она есть; правда, такая семантика
интерпретирует $\mathcal{U}_\ell$ как $\ell$-й
недостижимый кардинал.

\section{О чистой системе типов, основанной на рангах}

(TODO: завершить формализацию)

\begin{definition}
\emph{Рангом} функции называется максимальная
вложенность $\Pi$ в позиции аргумента:
\[\mathrm{rk}\,(\Pi(x:T).U)=\max\,(\mathrm{rk}\,T+1,
    \mathrm{rk}\,U)\]
\end{definition}

Рассмотрим $S=\{\mathrm{rk}_n\;|\;n\in\mathbb{N}\}$, и
\[\mathrm{Rule}=\{(\mathrm{rk}_m,\mathrm{rk}_n,
    \mathrm{rk}_{\max(m+1,n)})\;|\;m,n\in\mathbb{N}\}.\]
В качестве $\mathrm{Ax}$ можно взять цепочку
$\mathrm{rk}_0:\mathrm{rk}_1:\mathrm{rk}_2:\ldots$, но
получится система, строго менее удобная, чем обычная
иерархия вселенных. Попробуем поступить следующим
образом:
\[\mathrm{Ax}=\{(\mathrm{rk}_n,\mathrm{rk}_0)\;|\;
    n\in\mathbb{N}\}\]
В частности, $\mathrm{rk}_0:\mathrm{rk}_0$. Но в данной
системе парадокс Жирара не применим!

\subsection{Сильная нормализация}

Мы доказываем сильную нормализуемость с помощью аргумента кандидатов
приводимости, используя тот факт, что система \textbf{предикативна}: правило
произведения строго повышает уровень сорта, поэтому, хотя
\(\mathsf{Sort}_0:\mathsf{Sort}_0\), сорт \(\mathsf{Sort}_0\) не замкнут
относительно произведений.

\subsubsection{Кандидаты приводимости}

\begin{definition}
\textbf{Кандидат} — это множество \(C\subseteq \mathrm{SN}\), такое что:
\begin{enumerate}
\item Если \(t\in C\) и \(t\to t'\), то \(t'\in C\).
\item Если \(t\) нейтрален и всякий его непосредственный редукт лежит в \(C\), то \(t\in C\).
\end{enumerate}
\end{definition}

Множество \(\mathrm{SN}\) всех сильно нормализуемых термов само является кандидатом.

Определим кандидаты \(C_n\) для каждого сорта \(\mathsf{Sort}_n\) индукцией по \(n\).

\begin{itemize}
\item \(C_0 := \mathrm{SN}\).
\item Для \(n>0\): \(C_n\) — наименьший кандидат, содержащий:
\begin{itemize}
\item всякую переменную;
\item всякое произведение \(\Pi x{:}A.B\), такое что
\[
A\in C_m,\qquad 
\forall u\in \llbracket A\rrbracket,\ B[u]\in C_k,
\qquad \max(m+1,k)=n.
\]
\end{itemize}
\end{itemize}

Здесь \(\llbracket A\rrbracket\) — кандидат, интерпретирующий тип \(A\),
определённый ниже.

Это определение обосновано, поскольку в произведении сорта \(n\) область имеет
сорт \(m<n\), так что \(C_m\) уже определено. Кодомен может иметь сорт
\(k\le n\), но он встречается как подтерм, поэтому рекурсия структурна по
термам.

Затем доказывается индукцией по структуре термов, что каждое \(C_n\)
действительно является кандидатом.

\subsubsection{Интерпретация типов}

Для окружения \(\rho\), отображающего переменные типов в кандидаты, а термовые
переменные — в элементы соответствующих кандидатов, определим:
\[
\llbracket \mathsf{Sort}_n \rrbracket_\rho := C_n,
\]
\[
\llbracket x \rrbracket_\rho := \rho(x),
\]
\[
\llbracket \Pi x{:}A.B \rrbracket_\rho
:=
\Bigl\{ t \mid t\in\mathrm{SN}\ \wedge\
\forall u\in \llbracket A\rrbracket_\rho,\ t\,u\in \llbracket B\rrbracket_{\rho[x:=u]} \Bigr\}.
\]

Для типа \(A\) сорта \(\mathsf{Sort}_m\) индукцией по \(A\) показывается, что
\(\llbracket A\rrbracket_\rho\) — кандидат и что \(A\in C_m\).

Ключевой момент — случай произведения: если \(A\in C_m\) и для всех
\(u\in\llbracket A\rrbracket\) выполнено \(B[u]\in C_k\), то по определению
\(C_{\max(m+1,k)}\) произведение \(\Pi x{:}A.B\) лежит в \(C_{\max(m+1,k)}\).

\subsubsection{Основная лемма}

\begin{lemma}
Если \(\Gamma\vdash t:A\), то для всякого окружения \(\rho\), удовлетворяющего
\(\Gamma\), имеем \(t\in \llbracket A\rrbracket_\rho\).
\end{lemma}

\begin{proof}
Индукцией по выводу \(\Gamma\vdash t:A\).
\begin{itemize}
\item \textbf{Переменная:} непосредственно из того, что \(\rho\) удовлетворяет \(\Gamma\).
\item \textbf{Сорт:} \(\mathsf{Sort}_n:\mathsf{Sort}_0\). Нужно
    \(\mathsf{Sort}_n\in C_0=\mathrm{SN}\), что выполнено, поскольку сорта
        сильно нормализуемы.
\item \textbf{Произведение:} Если \(\Gamma\vdash A:\mathsf{Sort}_m\) и
    \(\Gamma,x:A\vdash B:\mathsf{Sort}_n\), то
    \(\Pi x{:}A.B:\mathsf{Sort}_{\max(m+1,n)}\). По индукции \(A\in C_m\) и для
    всех \(u\in\llbracket A\rrbracket\) имеем \(B[u]\in C_n\). Следовательно,
    \(\Pi x{:}A.B\in C_{\max(m+1,n)}\) по определению \(C_{\max(m+1,n)}\).
\item \textbf{Применение:} Если \(\Gamma\vdash f:\Pi x{:}A.B\) и
    \(\Gamma\vdash u:A\), то \(f\,u:B[u]\). По индукции
    \(f\in\llbracket \Pi x{:}A.B\rrbracket\) и \(u\in\llbracket A\rrbracket\).
    По определению интерпретации \(\Pi\) получаем
    \(f\,u\in\llbracket B[u]\rrbracket\).
\item \textbf{Абстракция:} Если \(\Gamma,x:A\vdash t:B\) и
    \(\Gamma\vdash \Pi x{:}A.B:\mathsf{Sort}_m\), то
    \(\lambda x{:}A.t:\Pi x{:}A.B\). Чтобы показать
    \(\lambda x{:}A.t\in\llbracket \Pi x{:}A.B\rrbracket\), возьмём произвольное
    \(u\in\llbracket A\rrbracket\). Тогда
\[
(\lambda x{:}A.t)\,u \to t[x:=u].
\]
По индукции \(t[x:=u]\in\llbracket B[u]\rrbracket\), а кандидаты замкнуты
относительно редукции, поэтому
\((\lambda x{:}A.t)\,u\in\llbracket B[u]\rrbracket\). Кроме того,
\(\lambda x{:}A.t\) сильно нормализуем, так как лежит в кандидате.
\end{itemize}
Таким образом, всякий хорошо типизированный терм лежит в некотором кандидате, а
значит, и в \(\mathrm{SN}\).
\end{proof}
